% Generated by scripts/assemble.py; corrections belong in source/corrections.json.
\examhead{196}{湘潭大学 2024 年数学分析试卷}{185}
% p186-b10; source PDF page 186
\par\Needspace{5\baselineskip}
\blockanchor{p186-b10}
\majorquestion{一}计算题. 每题 10 分, 共 30 分.\par
% p186-b11; source PDF page 186
\blockanchor{p186-b11}
(1) 求极限  $\lim_{x\to0}\frac{\tan x-\sin x}{x^{3}}$ .\par
% p186-b12; source PDF page 186
\blockanchor{p186-b12}
(2) 求极限 $\lim_{n\to \infty}\int_0^1\frac{x^n}{1 + x}\mathrm{d}x.$\par
% p186-b13; source PDF page 186
\blockanchor{p186-b13}
(3) 求极限  $\lim_{x\to+\infty}x^{\frac{1}{x}}$ .\par
% p186-b14; source PDF page 186
\par\Needspace{5\baselineskip}
\blockanchor{p186-b14}
\majorquestion{二}计算题. 每题 10 分, 共 20 分.\par
% p186-b15; source PDF page 186
\blockanchor{p186-b15}
(1) 设 $y = \frac{\sqrt{1 + x^2}}{x}$, 求 $\mathrm{d}^2 y$.\par
% p186-b16; source PDF page 186
\blockanchor{p186-b16}
(2) 设 $I(y) = \int_{y}^{y^2} \frac{\cos xy}{x} \, \mathrm{d}x$, 求 $I'(y)$.\par
% p186-b17; source PDF page 186
\par\Needspace{5\baselineskip}
\par\Needspace{6\baselineskip}
\blockanchor{p186-b17}
\majorquestion{三}(15 分) 设 $\Omega$ 为球面 $x^{2} + y^{2} + z^{2} = 3$ 和抛物面 $2z = x^{2} + y^{2}$ 所围成的区域, 计算重积分\par
\nopagebreak[4]
% p186-b18; source PDF page 186
\blockanchor{p186-b18}
\displayformula{\iiint_ {\Omega} (x + y + z) ^ {2} \mathrm{d} x \mathrm{d} y \mathrm{d} z.}
% p186-b19; source PDF page 186
\par\Needspace{5\baselineskip}
\blockanchor{p186-b19}
\majorquestion{四}(15 分) 已知 $\Sigma$ 为封闭曲面, $l$ 为固定向量, $n$ 为 $\Sigma$ 的外法向量, 证明: $\iint_{\Sigma} \cos(n, l) \, \mathrm{d}S = 0$.\par
% p186-b20; source PDF page 186
\par\Needspace{5\baselineskip}
\blockanchor{p186-b20}
\majorquestion{五}（20分）已知含参量积分 $I = \int_0^{+\infty}e^{-\alpha x}\sin x\mathrm{d}x.$\par
% p187-b00; source PDF page 187
\blockanchor{p187-b00}
(1) 对于 $\alpha \geq \alpha_0 > 0$, 讨论 $I$ 的一致收敛性.\par
% p187-b01; source PDF page 187
\blockanchor{p187-b01}
(2) 对于 $\alpha > 0$, 讨论 $I$ 的一致收敛性.\par
% p187-b02; source PDF page 187
\par\Needspace{5\baselineskip}
\blockanchor{p187-b02}
\majorquestion{六}(15 分) 叙述 $\lim_{x \to x_0^+} f(x)$ 存在且有限的柯西收敛准则并证明.\par
% p187-b03; source PDF page 187
\par\Needspace{5\baselineskip}
\blockanchor{p187-b03}
\majorquestion{七}(15 分) 将 $f(x) = x^2$ 在 $[0, 2\pi]$ 上傅里叶展开.\par
% p187-b04; source PDF page 187
\par\Needspace{5\baselineskip}
\blockanchor{p187-b04}
\majorquestion{八}(20 分) 设 $\zeta(s) = \sum_{n=1}^{\infty} \frac{1}{n^s}$, $E(s) = \prod_{n=1}^{\infty} \frac{1}{1 - p_n^s}$, $s > 1$, 其中 $p_n$ 为第 $n$ 个质数.\par
% p187-b05; source PDF page 187
\blockanchor{p187-b05}
(1) 证明: $\zeta(s)$ 与 $E(s)$ 均收敛.\par
% p187-b06; source PDF page 187
\blockanchor{p187-b06}
(2) 证明: $\zeta(s) = E(s)$.\par
